\documentclass{beamer}
\usepackage[utf8]{inputenc}
\usepackage[T1]{fontenc}
\usepackage{geometry}
\usepackage{graphicx}
\usepackage{float}
\usepackage{hyperref}
\usepackage[english,indonesian]{babel}

\title{Optimasi Pemrosesan Template LaTeX Menggunakan Pendekatan Berbasis AI}
\author{Academic Support Agent}
\date{\today}

\begin{document}
\maketitle

\begin{abstract}

\end{abstract}

\chapter{PENDAHULUAN}


\chapter{TINJAUAN PUSTAKA}


% Original content (commented out):
% %%%%%%%% NO MODIFICAR NADA EN ESTE ARCHIVO

\documentclass{beamer}

\input{archivos/librerias}

% AGREGAR LOS DATOS DEL USUARIO
\input{datos.tex}

% LOGO DE LA FACULTAD
\newcommand{\objeto}{logos/logo_\facultad.png}
\newcommand{\colour}{color\facultad}
\newcommand{\urlfacultad}{www.\facultad.espol.edu.ec}


\input{archivos/ambientes}

\input{archivos/formato}

%COLOR DEL CAPTION
\DeclareCaptionFont{colour}{\color{\colour}}
\captionsetup{labelfont = colour}

\begin{document}
\begin{frame} %Página de inicio
  \titlepage
\end{frame}

\ifthenelse{\equal{\secciones}{Y}}
{
    \begin{frame} % EN CASO DE QUE SE UTILICEN SECCIONES, GENERAR TABLA DE CONTENIDO
        \frametitle{Contenidos}
        \tableofcontents
    \end{frame}
}

\input{cuerpo}

\end{document}


\end{document}